
In Section \ref{sec:deform-cone}, we set up notation and reviewed
known results.  Given this information, we are ready to begin the
proof of Theorem \ref{thm:main-cusp}.

\subsection{Boundary relations}

Restrict to the case when we have a single cusp.

We will let $U$ be a horoball neighborhood of the one cusp, and $V_1$,
\dots, $V_n$ be tubular neighborhoods of the components of the
singular locus $\Sigma$, each with radius $R$.  Here, we are letting
$n$ be the total number of components of $\Sigma$.  We will assume
throughout that these neighborhoods are chosen such that all
intersections of distinct neighborhoods are trivial.

%% Aug 6, 2007:  Set up omega, etc, and recall time dependency.
Let $\omega \in H^1(X;E)$ correspond to a harmonic infinitesimal
deformation of hyperbolic structure at time $t$, and decompose
$\omega$ in tubular neighborhoods and horoball neighborhoods as in the
previous section, and consider boundary terms.
%%

By equation (\ref{eqn:sum>0}), we see immediately that
\begin{equation}
-b_U(\eta_0, \eta_0) \leq \sum_{j=1}^n b_j.
\label{eqn:eta-bound}
\end{equation}
Again, note that $\eta_0$ in the term $b_U(\eta_0,\eta_0)$ refers to
the decomposition of $\eta$ valid only in the horoball neighborhood
$U$.
%% Aug 6, 2007
Also, recall that the terms of equation (\ref{eqn:eta-bound}) all
depend on time $t$, but we are suppressing $t$ for notational
purposes.
%%


Now, we will analyze the left hand side of (\ref{eqn:eta-bound}).
Recall that in a horoball neighborhood of a cusp, $\omega_0$ can be
written as a complex multiple of the form $\omega_\ell =
\omega_\ell(\infty)$ in equation (\ref{eqn:omega_l-cusp}).  That is,
there are real numbers $a$ and $b$
%% Aug 7, 2007
(for fixed $t$)
%%
such that
\begin{equation}
\omega_0 = (a+ib)\left[
  \begin{array}{ccc}
    0&0&0\\
    0&-1&-i\\
    0&-i&1
  \end{array}\right]
\label{eqn:eta_0-matrix}
\end{equation}

\begin{lemma}
  Using the notation above,
  $$-b_U(\eta_0, \eta_0) = 2(a^2+b^2)\mbox{Area}(\partial U).$$
  \label{lemma:b_U=area}
\end{lemma}

\begin{proof}
  By equation (\ref{eqn:eta_0-matrix}), we may write $\eta_0 =
  \mbox{Re}(\omega_0)$ and $*D\eta_0 = \mbox{Im}(\omega_0)$ as:
\begin{equation}
  \eta_0 = \left[\begin{array}{ccc}
      0&0&0\\
      0&-a&b\\
      0&b&a\end{array}\right] =
  \left[\begin{array}{c}0\\-a\\b\end{array}\right] \omega_2 +
  \left[\begin{array}{c}0\\b\\a\end{array}\right] \omega_3
\label{eqn:eta0mat}
\end{equation}
and:
$$
  *D\eta_0 = \left[\begin{array}{ccc}
      0&0&0\\
      0&-b&-a\\
      0&-a&b\end{array}\right] =
  \left[\begin{array}{c}0\\-b\\a\end{array}\right] \omega_2 +
  \left[\begin{array}{c}0\\-a\\b\end{array}\right] \omega_3.
$$
Here $\omega_1$, $\omega_2$, and $\omega_3$ are forms dual to the
vectors $e_1$, $e_2$, and $e_3$ giving cylindrical coordinates on
$U$.  Recall in particular that $e_1$ is radial, and $e_2$ is
tangent to a meridian.

We then may compute $-b_U(\eta_0, \eta_0)$ explicitly in terms of the
constants $a$ and $b$.
$$-b_U(\eta_0, \eta_0) = \int_{\partial U} \eta_0\wedge *D\eta_0 =
\int_{\partial U} 2(a^2+b^2)\omega_2\wedge\omega_3.$$

Since $\omega_2\wedge\omega_3$ is the area form for the torus
$\partial U$, we have:
$$-b_U(\eta_0, \eta_0) = 2 (a^2+b^2)\mbox{Area}(\partial U).$$
\end{proof}

Notice that equation
(\ref{eqn:eta-bound}) and Lemma \ref{lemma:b_U=area} imply:
\begin{equation}
0 \leq \sum_j b_j,
\label{eqn:b_j>0}
\end{equation}
and that we have equality by equation (\ref{eqn:sumbj=eta}) if and
only if $\eta$ is trivial on $X-V$, which happens only when the
deformation is trivial.


\subsection{Teichm{\"u}ller Space}
We now turn our attention to the term $(a^2+b^2)$ in Lemma
\ref{lemma:b_U=area}.  This term is closely related to the change
in metric of the torus $\partial U$.

Under the hyperbolic metric on $X=M-\Sigma$, the boundary $\partial U$
inherits a Euclidean structure.  This gives a point in the
Teichm\"uller space of the torus, $\mathcal{T}(T^2)$.  

\begin{remark}
We view $\mathcal{T}(T^2)$ as the space of flat structures on the
torus.  To each flat structure $\zeta$ is associated a unit area
Euclidean metric $g_{\zeta}$, and vice versa.  However, in the
following discussion we will refer to both the flat structure and its
associated metric, and so we will continue to distinguish between the
two.  We use Greek letters $\gamma$, $\zeta$ to denote points in
$\mathcal{T}(T^2)$, and $g$ with appropriate subscripts to denote the
induced metric on the torus.
\end{remark}

As the metric on $X$ changes under the cone deformation, the point in
$\mathcal{T}(T^2)$ will also change, tracing out a smooth path
$\gamma(t)$, with tangent vector $\gamma'(t)$.  The tangent vector
$\gamma'(t)$ is an infinitesimal change of the flat structure on the
torus.  Our 1--form $\omega$ encodes the infinitesimal deformation of
the hyperbolic cone structure $X_t$ on $X=M-\Sigma$.  So we may relate
$\gamma'(t)$ to $\omega$, or more particularly, to $a^2+b^2$.  We will
do so using the Teichm\"uller metric on $\mathcal{T}(T^2)$.

Definitions of the Teichm\"uller distance vary by constant factors in
the literature.  We will use the convention that the Teichm\"uller
distance between two points, $\sigma$ and $\tau$ in
$\mathcal{T}(T^2)$, is defined to be
$$d(\sigma, \tau) = \frac{1}{2}\inf_{f}\;\log K_f,$$ where $f\co
\sigma \to \tau$ is a $K_f$--quasiconformal map, with $K_f$ the
smallest such constant.  (See for example \cite{lehto}.)

%% \begin{remark}
%% For background on the Teichm\"uller space of the torus, including a
%% proof that it is isometric to the hyperbolic plane $\HH^2$, we refer
%% the reader to \cite{lehto}, or other references.  Note that the above
%% definition of the Teichm\"uller distance differs from that of
%% \cite{lehto} by a factor of $1/2$.  
%% \end{remark}

\begin{lemma}
  Let $\gamma(t)$ denote the path of $\partial U$ through
  $\mathcal{T}(T^2)$.  Let $\gamma'(t)$ denote its tangent vector.
  Then
  $$\sqrt{a^2+b^2} = \left\| \gamma'(t) \right\|,$$
  where the metric is the {T}eichm{\"u}ller metric.
\label{lemma:a+b-gamma}
\end{lemma}

\begin{proof}
Fix $t=t_0$, and consider the flat structure on the torus given by
$\partial U_{t_0} = \gamma(t_0)$.  We may assume it has unit area.
Recall that any Teichm\"uller geodesic through $\gamma(t_0)$ is given
as follows.  Select two orthogonal geodesic foliations $F_1$ and $F_2$
on the torus.  Select $\lambda>0$.  For any $u \in \RR$, obtain a new
metric by multiplying vectors along $F_1$ by $\exp(-\lambda\,u)$ and
multiplying vectors along $F_2$ by $\exp(\lambda\,u)$.  This is a
``stretch--squeeze'' map, which preserves the area of the torus.  It
gives a one--parameter family of flat structures on the torus, which
we denote by $\zeta(u)$.

Teichm\"uller's theorem says that in the case of the torus, this
stretch--squeeze map has the minimal quasiconformal distortion, so is
a Teichm\"uller geodesic.  For each $u$, the derivative of the map
$\zeta(u)$ takes an infinitesimal circle to an infinitesimal ellipse,
the ratio of whose axes is $\exp(2u\lambda)$.  Thus for any $u$,
the distance between the structure $\zeta(0) = \gamma(t_0)$ at time
$0$ and $\zeta(u)$ at time $u$ is
$$d(\zeta(0), \zeta(u)) = \frac{1}{2}\inf_f\log K_f =
\frac{1}{2}\log(e^{2u\lambda}) = \lambda\,u.$$ 
Therefore, $\lambda = \| \zeta'(u) \|$, where the norm is given with
respect to the Teichm\"uller metric.  In particular, $\lambda =
\|\zeta'(0)\|$.

On the other hand, for any point on the torus, let $u_1$ and $u_2$ be
orthonormal vectors in the directions of $F_1$ and $F_2$ respectively.
Then when we write the infinitesimal change of metric induced by
$\zeta'(0)$ in coordinates given by $u_1$ and $u_2$, we obtain at
every point a diagonal matrix $\Lambda$ with $-\lambda$, $\lambda$ on
the diagonal, since at each time $u$ those vectors are just multiplied
by $e^{-\lambda\,u}$, $e^{\lambda\,u}$, respectively.  Thus if we let
$g_u$ be the metric on the torus induced by $\zeta(u)$, we have
$$\left.\frac{d}{du}\right|_{u=0} g_u(x,y) = 2 g_0(\Lambda\,x, y).$$

Now we relate this discussion to the term $a^2+b^2$.

Recall the infinitesimal deformation of the hyperbolic cone structure
of $X$ is encoded by the one--form $\omega$.  In particular, the
infinitesimal change in metric is given by the real part $\eta$.  That
is, if we let $g_t$ be the metric on $X$ at time $t$, then 
\begin{equation}
\left.\frac{d}{dt}\right|_{t_0}\, g_t(x,y) \;=\; 2 g_{t_0}(\eta(t_0) x, y).
\label{eqn:dotg}
\end{equation}
(See the displayed equation on page 374 of \cite{hk:univ}.)

We are interested in the infinitesimal change in the metric on
$\partial U_{t_0}$.  Since $\eta$ decomposes into $\eta_0 + \eta_c$,
and only $\eta_0$ changes the holonomy of $\partial U$, we know the
infinitesimal change in metric on $\partial U_{t_0}$ is given by
$\eta_0$.

We have an explicit formula for $\eta_0$, written in orthonormal
coordinates with respect to the metric $g_{t_0}$, given by $e_1$,
$e_2$, $e_3$, where $e_1$ is radial and $e_2$ points in the direction
of the meridian.  Consider again equation (\ref{eqn:eta0mat}).  Since
the first row and column for the matrix of $\eta_0$ are zero, we see
that $\eta_0$ actually has no effect on or contribution to the change
of metric in the radial direction of $U$.  Hence $\eta_0$ itself gives
the infinitesimal change of metric on the torus $\partial U_{t_0}$, or
$\gamma'(t_0) = \eta_0(t_0)$.  

When we write $\eta_0$ in orthonormal coordinates $e_2$, $e_3$ on the
torus, we discard the first row and column of the matrix of equation
(\ref{eqn:eta0mat}), and obtain the $2\times 2$ matrix
\begin{equation}
A = \left[\begin{array}{cc}-a&b\\b&a\end{array}\right].
\label{eqn:define-A}
\end{equation}
Here the
first column of $A$ gives the infinitesimal change in the meridional
direction, the second the change in the direction orthogonal to the
meridian.  That is, if we restrict the metric $g_t$ to
$\partial U_t$, we may rewrite equation (\ref{eqn:dotg}) as
$$\left.\frac{d}{dt}\right|_{t=t_0}\, g_t(x,y) =
2g_{t_0}(\eta_0(t_0)\,x, y) = 2\left<Ax, y\right>.$$

Hence $A$ is a matrix representation of $\gamma'(t_0)$, the
infinitesimal change of metric of the torus $\partial U_{t_0}$.

We can diagonalize $A$.  For any point on the torus, there exist
orthonormal vectors $u_1$ and $u_2$ such that when we put $A$ into
coordinates given by $u_1$ and $u_2$, $A$ is diagonal with
$-\sqrt{a^2+b^2}$ and $\sqrt{a^2+b^2}$ on the diagonal.  Moreover,
since $A$ does not depend on the point on the torus (that is, for any
point on the torus we obtain the same matrix $A$), $u_1$ and $u_2$
determine orthogonal geodesic foliations $F_1$ and $F_2$.  These, in
turn, determine a Teichm\"uller geodesic $\zeta(u)$ with $\zeta(0) =
\gamma(t_0)$, whose initial tangent vector $\zeta'(0)$ agrees with $A$
at every point, and has Teichm\"uller norm $\|\zeta'(0\| =
\sqrt{a^2+b^2}$.  Thus infinitesimally, $\gamma(t)$ and $\zeta(u)$
agree near $t=t_0$ and $u=0$, respectively.  Hence $\|\gamma'(t_0)\| =
\|\zeta'(0)\| = \sqrt{a^2+b^2}$.
\end{proof}

\begin{remark}
In the proof of Lemma \ref{lemma:a+b-gamma}, we showed that for any
time $t$, $\gamma(t)$ agrees infinitesimally with an explicit
Teichm\"uller geodesic which is completely determined by the matrix
$A$.  We will use this again below.
\end{remark}

Finally, putting Lemmas \ref{lemma:b_U=area}, and
\ref{lemma:a+b-gamma} together with equation (\ref{eqn:eta-bound}), we
have completed the proof of the following theorem.

\begin{theorem}
	Let $X_t$ be a hyperbolic cone deformation with a cusp.  That is,
  for each time $t$, the hyperbolic cone manifold $X_t$ has a cusp
  which remains a cusp throughout the deformation.  Let $U_t$ be a
  horoball neighborhood of the cusp.  Let $\gamma(t)$ denote the path
  $\partial U_t$ travels through the Teichm{\"u}ller space of the
  torus.  Then the change in $\partial U_t$ is bounded by the following
  inequality.
$$2\,\|\gamma'(t)\|^2 \mbox{Area}(\partial U_t) \leq \sum b_j(t)$$
The sum on the right hand side is over all components of the singular
locus.  
\label{thm:cusp-1}
\end{theorem}

Note Theorem \ref{thm:cusp-1} gives an inequality identical to that of
Theorem \ref{thm:main-cusp}, except for the boundary terms $b_j$ on
the right hand side.  We will bound these, and complete the proof of
Theorem \ref{thm:main-cusp}, in Section \ref{sec:bdrybound}.  First,
we set up a similar theorem to Theorem \ref{thm:main-norm-length}.


\subsection{Normalized lengths on the torus}
	\label{sec:normlength}

We are interested in determining how \emph{normalized lengths} of
curves on the torus $\partial U$ change over the deformation. 

\begin{define}
  Let $\beta$ be a slope on the torus $T$, that is, an isotopy class
  of simple closed curves.  The normalized length of $\beta$ is
  defined to be:
  $$\mbox{Normalized length}(\beta) =
  \frac{\mbox{Length}(\beta)}{\sqrt{\mbox{Area}(T)}},$$
	where $\mbox{Length}(\beta)$ is defined to be the length of a
  geodesic representative of $\beta$ on $T$.
\end{define}

\begin{lemma}
	Let $\beta$ be a slope on $\partial U_t$.  Let $L_t(\beta)$ denote
	its normalized length, with derivative $\dot{L}(\beta)$. Then
	$L_t(\beta)$ satisfies:
\begin{equation}
  -\|\gamma'(t)\| \, L_t(\beta) \;\leq\; \dot{L}(\beta) \;\leq\;
   \|\gamma'(t)\| \,L_t(\beta).
  \label{eqn:norm-length-L}
\end{equation}
\label{lemma:ldot-bound}
\end{lemma}

\begin{proof}
Fix a time $t=t_0$.  We will show the lemma for $t_0$.

We saw in the proof of Lemma \ref{lemma:a+b-gamma} that $\gamma(t)$
agrees infinitesimally with an explicit Teichm\"uller geodesic
$\zeta(u)$ determined by the matrix $A$ of equation
(\ref{eqn:define-A}) when $t=t_0$ and $u=0$.  More specifically,
$\zeta(u)$ multiplies vectors along $F_1$ by
$\exp(-u\sqrt{a^2+b^2})$, and multiplies vectors along $F_2$ by
$\exp(u\sqrt{a^2+b^2})$, where $F_1$ and $F_2$ are orthogonal
geodesic foliations (determined by the matrix $A$) on the torus
$\partial U_{t_0}$.  Thus we will prove the lemma by showing it is
true when the change in metric on $\partial U_{t_0}$ is given by
$\zeta(u)$.

Now, $L_t(\beta)$ is defined to be the length of a geodesic
representative of $\beta$ divided by $\sqrt{\rm{Area}(\partial U_t)}$.
Without loss of generality, we may assume that $\rm{Area}(\partial
U_t) = 1$, since $\gamma'(t)$ is given by a trace free (hence area
preserving) matrix $A$.  Thus $L_t(\beta)$ is just the length of a
geodesic representative of $\beta$ at time $t_0$.

Let $\hat{\beta}_{t_0}\co I \to \partial U_{t_0}$ be a geodesic
representative of $\beta$ at time $t_0$.  Then $\hat{\beta}_{t_0}$
makes some angle $\theta$ with the foliation $F_1$.  For any $u$,
$\zeta(u)$ takes $\hat{\beta}_{t_0}$ to a new geodesic, still with
slope $\beta$, and length $L_u(\beta) = \alpha_u(\theta) \,
L_{t_0}(\beta)$.  Here
$$\alpha_u(\theta)
= \sqrt{\cos^2\theta \,\exp(-2u\sqrt{a^2+b^2}) +\sin^2\theta \exp(2u\sqrt{a^2+b^2})}$$
is maximized when $\theta=\pi/2$ and minimized when $\theta=0$.  Thus
we have
$$\exp(-\sqrt{a^2+b^2}\,u)\,L_{t_0}(\beta) \leq L_u(\beta) \leq
\exp(\sqrt{a^2+b^2}\,u) \,L_{t_0}(\beta).$$

Since $\left.\frac{d}{dt}\right|_{t=t_0} L_t(\beta) =
\left.\frac{d}{du}\right|_{u=0} L_u(\beta),$
$$-\sqrt{a^2+b^2} \, L_{t_0}(\beta) \leq \left.\frac{d}{dt}\right|_{t=t_0}
L_t(\beta) \leq \sqrt{a^2+b^2}\, L_{t_0}(\beta).$$
\end{proof}

Combining this with the inequality of Theorem \ref{thm:cusp-1},
we obtain the following theorem.

\begin{theorem}
  Let $L_t$ be the normalized length of a slope on a cusp torus
  $\partial U_t$ under a cone deformation, with derivative $\dot{L} =
  \frac{d}{dt} L_t$.  Then the change in $L_t$ is bounded by the
  following inequality:
  $$
  -\left(\sqrt{\frac{\sum b_i(t)}{2\mbox{Area}(\partial
        U_t)}}\right) \, L_t
	\;\leq\; \dot{L} \;\leq\;
  \left(\sqrt{\frac{\sum b_i(t)}{2\mbox{Area}(\partial U_t)}}\right)\,
  L_t.
  $$
\label{thm:norm-length-1}
\end{theorem}

\begin{proof}
  Solving for $\|\gamma'(t)\|$ in the inequality of Theorem
  \ref{thm:cusp-1}, we find
  $$\|\gamma'(t)\| \leq \sqrt{\frac{\sum b_j(t)}{2\mbox{Area}(\partial
      U_t)}}.$$
  The result is given by substituting this into equation
  (\ref{eqn:norm-length-L}).
\end{proof}


\subsubsection{Normalized height}

For our applications, the curve on the torus we are particularly
interested in is the one running orthogonal to the meridian.  When we
are dealing with link complements in $S^3$, as in a future section of
this paper, there is a well defined meridian on a cusp torus.
Otherwise, choose the shortest nontrivial simple closed curve on the
initial cusp torus to be the meridian (for our applications to knots
and links in $S^3$, these choices will agree).

We will be interested in estimating the length of the next shortest
nontrivial simple closed curve.  This length will be at least as long
as the length of the arc orthogonal to the meridian.  

At time $t$, let $p_t$ denote the geodesic arc perpendicular to the
meridian, running from meridian to meridian on $\partial U_t$.  This
arc will generally not be a closed curve.  Let $h_t$ be the normalized
length of $p_t$, that is, $h_t =
\rm{Length}(p_t)/\sqrt{\rm{Area}(\partial U_t)}$.  We refer to $h_t$
as the \emph{normalized height} of the torus at time $t$.  Denote its
derivative at time $t$ by $\dot{h}$.

\begin{lemma}
The normalized height $h_t$ satisfies
$$ -\sqrt{a^2+b^2}\,h_t \leq \dot{h} \leq
\sqrt{a^2+b^2}\,h_t.$$
\label{lemma:hdot-bound}
\end{lemma}

\begin{proof}
%% Jan 11, 2008.  Modify according to referee's suggestion.
Let $\mu$ denote the meridian slope, with $L_t(\mu)$ its normalized
length at time $t$.  Then $h_t L_t(\mu) = 1$.  Thus
$$\frac{\dot{h}}{h_t} + \frac{\dot{L}(\mu)}{L_t(\mu)} = 0.$$
By Lemma \ref{lemma:ldot-bound} and Lemma \ref{lemma:a+b-gamma},
$$-\sqrt{a^2+b^2} \leq \frac{\dot{L}(\mu)}{L_t(\mu)} \leq
\sqrt{a^2+b^2}.$$ Hence
$ -\sqrt{a^2+b^2}\,h_t \leq \dot{h} \leq
\sqrt{a^2+b^2}\,h_t.$

%% Fix $t=t_0$.  %% As in the proof of Lemma \ref{lemma:a+b-gamma},
%% %% $\gamma'(t_0) = \zeta'(0)$ where $\zeta(u)$ is the Teichm\"uller
%% %% geodesic which at time $u$ multiplies vectors along $F_1$ by
%% %% $\exp(-\sqrt{a^2+b^2}\, u)$ and multiplies vectors along $F_2$ by
%% %% $\exp(\sqrt{a^2+b^2}\,u)$, where $F_1$ and $F_2$ are orthogonal
%% %% geodesic foliations of the torus, along directions determined by the
%% %% matrix $A$.  Therefore again we prove the lemma for the case when
%% %% the change in metric on $\partial U_{t_0}$ is given by $\zeta(u)$.

%% %% Additionally, again since $A$ is traceless, we may assume that
%% %% $\partial U_{t}$ has unit area.  Thus $h_t$ is the length $L_t(p_t)$
%% %% measured on the torus $\partial U_t$.
%% Consider the geodesic arc $p_{t_0}$.  We may follow the proof of Lemma
%% \ref{lemma:ldot-bound} exactly to show that the length $L_u(p_{t_0})$
%% satisfies:
%% \begin{equation}
%% -\sqrt{a^2+b^2}\,L_{t_0}(p_{t_0}) \leq \left.\frac{d}{dt}\right|_{u=t_0}
%% L_t(p_{t_0}) \leq \sqrt{a^2+b^2}\,L_{t_0}(p_{t_0}),
%% \label{eqn:hdot}
%% \end{equation}
%% where here $L_u(p_{t_0})$ is the length of the geodesic which
%% $p_{t_0}$ becomes under $\zeta(u)$.

%% This does not quite finish the proof.  Although $p_{t_0}$ remains
%% geodesic under $\zeta(u)$, it will not necessarily remain
%% perpendicular to the meridian.  Thus $L_u(p_{t_0})$ will not in
%% general equal $h_t$.  However, we have
%% $$ \left.\frac{d}{dt} \right|_{t=t_0} h_t =
%% \left.\frac{d}{dt}\right|_{t=t_0} L_t(p_t) =
%% \left.\frac{d}{dt}\right|_{t=t_0} L_t(p_{t_0}) +
%% \left.\frac{d}{dt}\right|_{t=t_0} L_{t_0}(p_t).$$

%% At time $t=t_0$, $p_{t_0}$ is the geodesic arc perpendicular to
%% the meridian.  For $t$ near $t_0$, $p_t$ will be a geodesic in a
%% direction near $p_{t_0}$.  Any such arc must have length on $\partial
%% U_{t_0}$ at least as long as that of the perpendicular $p_{t_0}$.
%% Hence $L_{t_0}(p_t)$ has a local minimum at $t=t_0$.  So the term
%% $\left.\frac{d}{dt}\right|_{t=t_0} L_{t_0}(p_t)$ is equal to zero.  

%% Then $\dot{h} = \left.\frac{d}{dt}\right|_{t=t_0}
%% L_t(p_{t_0})$, so equation (\ref{eqn:hdot}) gives the result.
\end{proof}

Again Lemma \ref{lemma:hdot-bound}, combined with the inequality of
Theorem \ref{thm:cusp-1}, gives
\begin{equation}
  -\left(\sqrt{\frac{\sum b_i}{2\mbox{Area}(\partial U_t)}}\right) \,
h_t \leq \dot{h} \leq
\left(\sqrt{\frac{\sum b_i}{2\mbox{Area}(\partial U_t)}}\right) \,
h_t
\label{eqn:norm-height}
\end{equation}

We simplify equation (\ref{eqn:norm-height}) by noting the area of
$U_t$ is given by the length of the meridian times the actual length
of the arc $p_t$.  Recall that the length of $p_t$ is
$\sqrt{\mbox{Area}(\partial U_t)}\,L_t(p_t) =
\sqrt{\mbox{Area}(\partial U_t)}h_t.$ So denoting the length of the
meridian of $\partial U_t$ by $m_t$, we have
$$\sqrt{\mbox{Area}(\partial U_t)} = m_t h_t.$$

Putting this into equation (\ref{eqn:norm-height}), we obtain:
\begin{theorem}
  Let $M$ be a hyperbolic 3--manifold with a cusp, and let $X$ be a
  hyperbolic 3--manifold which can be joined to $M$ by a smooth family
  of hyperbolic cone manifolds.  Let $h_t$ denote the normalized
  height at time $t$ of the cusp which deforms to the cusp of $M$, and
  let $m_t$ denote the length of its meridian.  Finally, denote by
  $\dot{h}$ the derivative of $h_t$.  Then
\begin{equation}
  -\frac{\sqrt{\sum b_i(t)}}{m_t\sqrt{2}} \leq \dot{h} \leq
\frac{\sqrt{\sum b_i(t)}}{m_t\sqrt{2}}
\label{eqn:norm-height-mer}
\end{equation}
\label{thm:norm-height-mer}
\end{theorem}








